\documentclass[10pt,twocolumn]{article}

%------------------------------------------------------------
% PACKAGES
%------------------------------------------------------------
\usepackage[margin=0.75in,top=1in,bottom=1in]{geometry}
\usepackage{graphicx}
\usepackage{amsmath}
\usepackage{amssymb}
\usepackage{booktabs}
\usepackage{array}
\usepackage{hyperref}
\usepackage{xcolor}
\usepackage{caption}
\usepackage{float}
\usepackage{enumitem}
\usepackage{algorithm}
\usepackage{algorithmic}
\usepackage{cite}
\usepackage{microtype}
\usepackage{multicol}
\usepackage{titlesec}
\usepackage{abstract}
\usepackage{parskip}
\usepackage{fancyhdr}
\usepackage{lipsum}
\usepackage{multirow}
\usepackage{tabularx}

%------------------------------------------------------------
% HYPERREF SETUP
%------------------------------------------------------------
\hypersetup{
    colorlinks=true,
    linkcolor=blue!70!black,
    citecolor=blue!70!black,
    urlcolor=blue!70!black
}

%------------------------------------------------------------
% SECTION FORMAT
%------------------------------------------------------------
\titleformat{\section}{\normalfont\large\bfseries\centering}{\Roman{section}.}{0.5em}{}
\titleformat{\subsection}{\normalfont\normalsize\bfseries\itshape}{\Alph{subsection}.}{0.5em}{}
\titleformat{\subsubsection}{\normalfont\small\bfseries}{}{0em}{}

%------------------------------------------------------------
% ABSTRACT FORMAT
%------------------------------------------------------------
\renewcommand{\abstractnamefont}{\normalfont\bfseries}
\renewcommand{\abstracttextfont}{\normalfont\small\itshape}

%------------------------------------------------------------
% HEADER / FOOTER
%------------------------------------------------------------
\pagestyle{fancy}
\fancyhf{}
\rhead{\small\thepage}
\lhead{\small\textit{WSH++: Communication-Aware and Energy-Efficient Workflow Scheduling on Heterogeneous Hadoop Clusters}}
\renewcommand{\headrulewidth}{0.4pt}

%------------------------------------------------------------
% BEGIN DOCUMENT
%------------------------------------------------------------
\begin{document}

%------------------------------------------------------------
% TITLE BLOCK (one-column spanning header)
%------------------------------------------------------------
\twocolumn[{%
\begin{center}

{\LARGE \bfseries Communication-Aware and Energy-Efficient Workflow Scheduling\\
on Heterogeneous Hadoop Clusters via PSO and Genetic Algorithms}\\[0.5cm]

{\large \bfseries Sriman Srinivasan$^{1}$ \quad Neil Chetty$^{1}$}\\[0.25cm]

{\normalsize $^{1}$Department of Computer Science and Engineering,
Indian Institute of Information Technology Design and Manufacturing (IIITDM), Kurnool, Andhra Pradesh, India}\\[0.1cm]

{\normalsize \texttt{\{123AD0011, 123AD0003\}@iiitdm.ac.in}}\\[0.15cm]

{\small \textit{Roll Numbers:}
\textbf{Sriman Srinivasan} --- 123AD0011 \quad|\quad
\textbf{Neil Chetty} --- 123AD0003}\\[0.3cm]

{\small \textit{Under the Guidance of}}\\[0.1cm]
{\normalsize \textbf{Dr.~Srinivas Naik}}\\
{\small Associate Professor, Department of Computer Science and Engineering,\\
IIITDM Kurnool, Andhra Pradesh, India}\\[0.5cm]

{\small April 2026}\\[0.6cm]

\end{center}

%------------------------------------------------------------
% ABSTRACT
%------------------------------------------------------------
\begin{onecolabstract}
\noindent
Scheduling scientific workflows on heterogeneous Hadoop clusters is an NP-Hard
optimisation problem that becomes more complex in distributed storage environments
where inter-node data transfer costs are non-trivial. This paper presents
\textbf{WSH++}, a comprehensive extension of the Workflow Scheduler on Hadoop (WSH)
proposed by Rahmani et al.\ (2025). The original WSH algorithm addresses
heterogeneous Hadoop scheduling through subcluster partitioning and a training-task
mechanism, but explicitly acknowledges three core limitations: (i) communication
costs are set to zero, making scheduling decisions unrealistic in distributed storage;
(ii) only one baseline (HEFT) is evaluated; and (iii) energy consumption is ignored.
WSH++ systematically addresses all three limitations through: a
\textbf{communication-aware cost model} that accounts for intra-cluster latency,
inter-cluster bandwidth penalties, and data-size-driven transfer times; a
\textbf{Particle Swarm Optimisation (PSO) scheduler} and a \textbf{Genetic
Algorithm (GA) scheduler} that optimise a joint makespan--communication--energy
objective; and an \textbf{energy-aware scheduler} based on a physics-motivated power
model. All five schedulers are evaluated across three scientific workflows
(Gene2Life, Avianflu\_small, Epigenomics) and five node configurations (2--13 nodes)
on a Docker-emulated heterogeneous Hadoop cluster with CPU pinning. PSO achieves the
best makespan in 8 of 15 workflow--node configurations, outperforming WSH by up to
6.8\%. WSH consistently outperforms HEFT by 17--39\% on medium-to-large workflows.
The energy-aware scheduler reduces predicted energy consumption by 60--80\% relative
to HEFT, confirming that multi-objective scheduling is fully feasible within the WSH
framework with no meaningful makespan regression. The complete implementation is
provided as an open-source Java 17 system with Docker-based cluster emulation.

\smallskip
\noindent\textbf{Keywords:} Hadoop scheduling, heterogeneous computing, Particle
Swarm Optimisation, Genetic Algorithm, energy-aware scheduling, communication cost
model, workflow management, Hi-WAY, HEFT.
\end{onecolabstract}
\vspace{0.8cm}
}]

%------------------------------------------------------------
% I. INTRODUCTION
%------------------------------------------------------------
\section{Introduction}

The volume of data generated by IoT devices, genomics platforms, astrophysics
instruments, and social networks now reaches petabyte scales per day. Processing
such data requires distributed workflow management systems capable of orchestrating
hundreds of interdependent jobs across large computing clusters. Apache
Hadoop~\cite{hadoop2023} has established itself as the dominant platform for
large-scale data processing, providing fault-tolerant distributed storage through
HDFS and cluster resource management through YARN. When the underlying cluster is
\emph{heterogeneous}---nodes differ in CPU speed, memory capacity, and network
bandwidth---scheduling becomes significantly harder: the execution time of a job is
no longer a property of the job alone but of the node to which it is assigned.

Scientific workflow management systems such as Hi-WAY~\cite{hiway2017} sit atop
Hadoop YARN and abstract individual tool invocations as directed acyclic graph (DAG)
nodes. Scheduling in Hi-WAY requires mapping each DAG node (Hi-WAY \emph{job}) to a
Hadoop container while respecting data-dependency edges and minimising total workflow
completion time, known as \emph{makespan}. This problem is NP-Hard even on
homogeneous processors~\cite{hanen1994}, and the introduction of heterogeneity and
communication delays makes it substantially harder in practice.

The \emph{Heterogeneous Earliest Finish Time} (HEFT) algorithm by Topcuoglu
et al.~\cite{heft2002} is the canonical greedy heuristic for this class of
problems. HEFT sorts jobs by upward rank---a path-based priority measure---and
greedily assigns each job to the processor achieving the earliest finish time,
with optional idle-time insertion. While simple and effective on synthetic
benchmarks, HEFT was not designed for Hadoop-specific features such as data
locality, HDFS block placement, or the existence of IO-intensive versus
compute-intensive workloads.

The \emph{Workflow Scheduler on Hadoop} (WSH), proposed by Rahmani
et al.~\cite{rahmani2025}, addresses these Hadoop-specific characteristics through
two innovations. First, it partitions the heterogeneous cluster into homogeneous
subclusters, within each of which execution time estimates are uniform. Second,
before the scheduling phase proper, it runs a training task on the first node of
every subcluster to obtain actual finish-time benchmarks. These benchmarks rank
subclusters per job and allow IO-intensive jobs (exhibiting nearly equal training
times across all subclusters) to be routed to weaker nodes, freeing powerful nodes
for compute-intensive work. Rahmani et al.\ report a 42\% reduction in Scheduling
Length Ratio (SLR), a 20\% reduction in makespan, and a 37\% improvement in speedup
over HEFT across three scientific workflows.

\noindent\textbf{Motivation and Gaps.}
Despite these advances, Rahmani et al.\ explicitly identify four limitations and
directions for future work:
\begin{enumerate}[leftmargin=1.3cm,label=\arabic*.,topsep=2pt,itemsep=1pt]
  \item Communication costs between nodes are set to zero ($c_{ij}=0$), which is
        unrealistic when dependent jobs reside on different nodes that must exchange
        large intermediate datasets over the network.
  \item Only HEFT is used as a comparison baseline; no metaheuristic global
        optimisers are evaluated.
  \item Energy consumption is neither modelled nor minimised, despite green cloud
        computing being an increasingly important concern.
  \item Evaluation is limited to three workflows on a single private cloud
        configuration.
\end{enumerate}

\noindent\textbf{Contributions.}
This paper makes the following novel contributions, each directly addressing one or
more of the gaps listed above:
\begin{enumerate}[leftmargin=1.3cm,label=\arabic*.,topsep=2pt,itemsep=1pt]
  \item \textbf{Communication-Aware Cost Model.} We derive a latency-plus-bandwidth
        edge-weight model that distinguishes same-node, same-cluster, and
        cross-cluster data transfers, and integrate it into both the upward rank and
        the EST computation.
  \item \textbf{PSO Scheduler.} A Particle Swarm Optimisation scheduler that
        represents scheduling decisions as continuous bias vectors and minimises a
        weighted joint objective over makespan, communication, and energy.
  \item \textbf{GA Scheduler.} A Genetic Algorithm scheduler with elitism,
        tournament selection, single-point crossover, and Gaussian mutation,
        operating on the same three-objective fitness function.
  \item \textbf{Energy-Aware Scheduler.} A physics-motivated power model and a
        HEFT-style greedy placement strategy that penalises high-power nodes,
        achieving 60--80\% predicted energy reduction with no makespan regression.
  \item \textbf{Reproducible Open-Source Benchmark.} A complete Java 17
        implementation with Docker CPU pinning provides a reproducible heterogeneous
        Hadoop cluster on a single server for 75 scheduler runs across three
        workflows and five node configurations.
\end{enumerate}

The remainder of this paper is organised as follows. Section~\ref{sec:related}
surveys related work. Section~\ref{sec:wsh_original} reviews the original WSH
algorithm and its limitations. Section~\ref{sec:wshpp} details the WSH++
extensions. Section~\ref{sec:evaluation} presents experimental results.
Section~\ref{sec:discussion} discusses findings and limitations.
Section~\ref{sec:conclusion} concludes.

%------------------------------------------------------------
% II. RELATED WORK
%------------------------------------------------------------
\section{Related Work}
\label{sec:related}

\subsection{Greedy Heuristics for DAG Scheduling}

\textbf{HEFT}~\cite{heft2002} remains the most cited greedy heuristic for
scheduling task graphs on heterogeneous processors. It computes upward rank for
each task---a weighted longest-path estimate incorporating average execution and
average communication costs---and assigns tasks in decreasing rank order to
the processor yielding the earliest finish time. The optional idle-time insertion
variant allows a task to be placed into an idle slot earlier on a processor's
timeline, reducing makespan at the cost of additional bookkeeping. HEFT's
$O(n^2 p)$ time complexity (for $n$ tasks and $p$ processors) makes it fast enough
for online use, but its greedy nature means it can be trapped by locally optimal
assignments that are globally suboptimal.

\textbf{MRWS}~\cite{mrws2016} adapts HEFT to the MapReduce paradigm by converting
scientific workflow jobs into MapReduce tasks. MRWS introduces idle-time insertion
specifically for task slots within a MapReduce job and uses task duplication to
improve data locality. However, MRWS treats all nodes as a flat pool and does not
differentiate IO-intensive from compute-intensive jobs, which limits its
effectiveness on heterogeneous clusters.

\subsection{Metaheuristic Schedulers}

\textbf{HEPGA}~\cite{hepga} combines HEFT initialisation, PSO local refinement, and
GA global exploration in a three-phase hybrid for cloud workflow scheduling. Tested
exclusively on CloudSim, HEPGA demonstrates that metaheuristic post-processing can
improve HEFT solutions by 5--12\% on synthetic benchmarks. Our work differs by
operating on a real Hadoop cluster (emulated via Docker) rather than a simulator.

\textbf{Meshkati and Safi-Esfahani}~\cite{pso_cloud} apply PSO and Artificial Bee
Colony (ABC) algorithms to energy-aware resource allocation in cloud computing,
demonstrating that swarm-based methods are effective when the objective includes an
energy term. Their power model motivates the energy formulation adopted in WSH++.

\textbf{Yang et al.}~\cite{yang2022} propose classification-based diverse workflow
scheduling for multi-cloud environments, optimising makespan, cost, and energy
simultaneously. Their classification approach for job types parallels WSH's
IO/compute differentiation, but operates on WorkflowSim rather than a real Hadoop
platform.

\subsection{Energy-Aware Scheduling}

\textbf{Wang et al.}~\cite{wang2024} study electricity-cost-aware multi-workflow
scheduling in heterogeneous cloud environments, jointly optimising monetary cost and
makespan under time-varying electricity pricing. Their formulation treats energy as
an economic objective; WSH++ treats it as a physical quantity with a direct power
model tied to node hardware characteristics.

\textbf{DRL-YARN}~\cite{xue2023} applies deep reinforcement learning to workflow
scheduling on YARN clusters, reducing average job latency through learned policies.
While DRL-YARN eliminates the need for an explicit objective function, it requires
extensive offline training, introduces significant scheduling overhead, and
demonstrates limited generalisation to unseen workflow structures---limitations that
greedy and metaheuristic approaches do not share.

\subsection{Hadoop-Specific Scheduling}

The original WSH paper~\cite{rahmani2025} is the most direct antecedent of this
work. Its subcluster model and training-task mechanism are unique among Hadoop
schedulers in that they use real execution data (not simulation) to inform
scheduling decisions, making them more robust to workload variability. WSH++ builds
directly on this foundation and extends it rather than replacing it, preserving the
advantages of the training-task approach while addressing the three core limitations.

\noindent\textbf{Summary.} None of the existing works simultaneously (a) operates
natively on Hadoop/Hi-WAY with a real training-task mechanism, (b) extends a
training-based scheduler with a communication-aware cost model, and (c) evaluates
PSO, GA, and energy-aware schedulers on the same Docker-emulated Hadoop cluster
under identical conditions. WSH++ fills this gap.

%------------------------------------------------------------
% III. ORIGINAL WSH: SUMMARY AND LIMITATIONS
%------------------------------------------------------------
\section{Original WSH: Summary and Limitations}
\label{sec:wsh_original}

\subsection{Problem Formulation}

Let $\mathcal{W} = (J, E)$ be a workflow DAG where $J = \{j_1, \ldots, j_n\}$ is
the set of Hi-WAY jobs and $E \subseteq J \times J$ is the set of precedence edges.
An edge $(j_i, j_k) \in E$ means $j_k$ cannot start until $j_i$ has completed and
its output data has been transferred to the node where $j_k$ will run.

The cluster $\mathcal{C} = \{C_1, \ldots, C_K\}$ is a set of homogeneous
subclusters. Within each subcluster $C_k$, all nodes have the same CPU speed and
memory, so the execution time of job $j_i$ is the same on every node in $C_k$.
Let $w(j_i, C_k)$ denote this common execution time. The scheduling problem is to
produce an assignment $\phi : J \rightarrow \bigcup_k C_k$ and a start-time vector
$\mathbf{s}$ such that all precedence constraints are satisfied and the makespan
$C_{\max} = \max_{j \in J} (\text{AFT}(j))$ is minimised.

\subsection{Stage 1: Job Prioritisation}

WSH sorts jobs in descending order of upward rank. With communication costs zeroed
(as in the original paper):
\begin{align}
    \text{rank}_u(j_i) &= \overline{w}_i + \max_{j_k \in \text{succ}(j_i)}
    \left[c_{ik} + \text{rank}_u(j_k)\right]
    \label{eq:orig_rank}
\end{align}
where $\overline{w}_i = \frac{1}{K}\sum_{k=1}^{K} w(j_i, C_k)$ is the average
execution cost across all subclusters and $c_{ik} = 0$ in the original formulation.
For the exit job $j_{\text{exit}}$:
$\text{rank}_u(j_{\text{exit}}) = \overline{w}_{\text{exit}}$.

\subsection{Stage 2: Job Allocation}

Before scheduling, a training task is executed on the first node of each subcluster.
The training finish times $\text{FT}[j_i, C_k]$ rank the subclusters for job $j_i$.
Each job is then assigned to the container minimising:
\begin{equation}
    \text{EFT}(j_i, c_t) = \text{EST}(j_i, c_t) + w(j_i, c_t)
    \label{eq:orig_eft}
\end{equation}
\begin{equation}
    \text{EST}(j_i, c_t) = \max\!\left\{\text{readyTime}(j_i),\ \text{avail}(c_t)\right\}
    \label{eq:orig_est}
\end{equation}
\begin{equation}
    \text{readyTime}(j_i) = \max_{j_k \in \text{pred}(j_i)} \text{AFT}(j_k)
    \label{eq:orig_ready}
\end{equation}

If the training times of a job across subclusters are nearly equal (within 15\%),
the job is classified as IO-intensive and routed to weaker subclusters, freeing
powerful nodes for compute-intensive work.

\subsection{Identified Limitations}

The original paper explicitly acknowledges the following limitations:

\begin{enumerate}[leftmargin=1.3cm,label=(\roman*),topsep=2pt,itemsep=1pt]
  \item \textbf{$c_{ij} = 0$ throughout.} Setting all edge communication costs to
        zero ignores the reality that dependent jobs on different nodes or subclusters
        must transfer potentially large intermediate datasets over the network. This
        makes EST computations optimistic and can lead to scheduling decisions that
        concentrate data-dependent jobs on nodes with high inter-cluster latency.
  \item \textbf{Only HEFT as baseline.} A single comparison point cannot establish
        whether WSH's training-task heuristic captures most of the achievable
        improvement or whether a global search could do significantly better.
  \item \textbf{No energy model.} With power consumption of data centres accounting
        for 1--2\% of global electricity usage~\cite{wang2024}, ignoring energy is a
        significant practical limitation.
  \item \textbf{Future work.} The authors explicitly call for PSO/GA integration and
        communication-cost inclusion as the two primary future directions.
\end{enumerate}

%------------------------------------------------------------
% IV. WSH++: EXTENDED SCHEDULING FRAMEWORK
%------------------------------------------------------------
\section{WSH++: Extended Scheduling Framework}
\label{sec:wshpp}

WSH++ retains the subcluster partitioning and training-task mechanism of the
original WSH and augments it with four interacting extensions: a
communication-aware cost model, an energy consumption model, a PSO scheduler, and a
GA scheduler, all sharing a unified three-term objective function.

\subsection{Communication-Aware Cost Model}

Let $\text{bytes}(p)$ denote the modelled output size (in megabytes) of parent job
$p$, $B_u$ and $B_v$ the network bandwidths of nodes $u$ and $v$ respectively, and
$L_s = 4$\,ms and $L_c = 18$\,ms the same-cluster and cross-cluster base latencies.
The edge communication cost for a data dependency from $p$ (allocated to node $u$)
to its successor $j$ (being considered for node $v$) is:

\begin{equation}
    c_{p,j}(u,v) = \begin{cases}
        0 & u = v \\[4pt]
        L_s + \dfrac{\text{bytes}(p) \cdot 10^3}{\min(B_u, B_v)} & \text{same cluster} \\[8pt]
        L_c + 3 \cdot \dfrac{\text{bytes}(p) \cdot 10^3}{\min(B_u, B_v)} & \text{cross cluster}
    \end{cases}
    \label{eq:comm_cost}
\end{equation}

The factor of 3 for cross-cluster transfers captures overhead from inter-switch
routing, TCP/IP framing, and HDFS block replication. The average communication cost
$\overline{c}_{ik}$ used in rank computation is the mean of $c_{p,j}(u,v)$ over all
node-pair combinations $(u,v)$.

\noindent\textbf{Extended upward rank:}
\begin{equation}
    \text{rank}_u(j_i) = \overline{w}_i + \max_{j_k \in \text{succ}(j_i)}
    \!\left[\overline{c}_{ik} + \text{rank}_u(j_k)\right]
    \label{eq:ext_rank}
\end{equation}

\noindent\textbf{Communication-aware EST:}
\begin{equation}
    \text{EST}(j_i, c_t) = \max\!\left[\text{avail}(c_t),\
    \max_{p \in \text{pred}(j_i)}\!\left(\text{AFT}(p) + c_{p,j_i}\right)\right]
    \label{eq:ext_est}
\end{equation}

This formulation makes EST depend not only on when a container becomes free but
also on when predecessor data physically arrives from the node where the predecessor
ran. Jobs with large input data from cross-cluster predecessors therefore incur a
meaningful delay that the original WSH silently ignores.

\textbf{Impact on scheduling decisions.}
With non-zero communication costs, the scheduler is naturally incentivised to
co-locate dependent jobs within the same subcluster---or even the same node---when
the execution-time penalty of doing so is less than the communication saving. This
produces qualitatively different schedules for workflows with large intermediate
datasets, such as Epigenomics, where FASTQ files transferred between the Split and
Filter stages can be hundreds of megabytes.

\subsection{Energy Consumption Model}

The predicted energy consumed by job $j_i$ running on node $v$ for $t_{\text{exec}}$
milliseconds is:
\begin{equation}
    E(j_i, v) = P_{a} \cdot \alpha_{j_i} \cdot t_{\text{exec}}
                + 0.1\,P_{a}\,t_{\text{exec}}
                + 6\,t_s + 14\,t_x
    \label{eq:energy}
\end{equation}
where:
\begin{itemize}[leftmargin=1.4cm,nosep]
  \item $P_{a}$ is the node's modelled active power (W), derived from its vCPU and
        memory configuration (C1: 80\,W, C2: 60\,W, C3: 45\,W, C4: 30\,W).
  \item $\alpha_{j_i} \in \{1.0,\ 0.82\}$ is the task-type load factor: 1.0 for
        compute-intensive jobs (full CPU utilisation) and 0.82 for IO-intensive
        jobs (memory/disk-bound, lower CPU load).
  \item $0.1\,P_{a}$ is the idle leakage power (10\% of active power).
  \item $t_s$, $t_x$ are same-cluster and cross-cluster transfer durations (s).
  \item 6\,W and 14\,W are network-interface energy rates for intra- and
        inter-cluster transfers respectively.
\end{itemize}

The total predicted energy for a schedule $\phi$ is:
\begin{equation}
    E_{\text{total}} = \sum_{j_i \in J} E(j_i, \phi(j_i))
    \label{eq:energy_total}
\end{equation}

\subsection{Joint Objective Function}

All metaheuristic schedulers minimise a three-term weighted objective:
\begin{equation}
    f(\phi) = C_{\max}(\phi)
              + \lambda_1 \cdot C_{\text{comm}}(\phi)
              + \lambda_2 \cdot E_{\text{total}}(\phi)
    \label{eq:obj}
\end{equation}
where $C_{\text{comm}}(\phi) = \sum_{(p,j)\in E} c_{p,j}(\phi(p),\phi(j))$ is the
total predicted communication time. After grid search over $\lambda_1 \in [0.005,
0.1]$ and $\lambda_2 \in [0.05, 0.5]$ on the Gene2Life workflow, we set
$\lambda_1 = 0.020$ and $\lambda_2 = 0.14$. These weights balance the three terms
such that neither communication nor energy dominates for typical workloads.

\subsection{PSO Scheduler}

\textbf{Encoding.}
Each particle represents a scheduling policy as a continuous triple:
\begin{equation}
    \mathbf{x} = (\mathbf{b}_{\text{prio}},\, \mathbf{b}_{\text{clus}},\,
                  \mathbf{b}_{\text{energy}}) \in \mathbb{R}^{3n}
\end{equation}
where $n$ is the number of jobs. The priority bias $\mathbf{b}_{\text{prio}}$
perturbs the upward-rank ordering; the cluster bias $\mathbf{b}_{\text{clus}}$
shifts subcluster preferences; and the energy bias $\mathbf{b}_{\text{energy}}$
modulates how aggressively the scheduler avoids high-power nodes. A concrete
schedule is constructed from each particle by (i) reordering jobs by
$\text{rank}_u(j_i) + b_{\text{prio},i}$, (ii) assigning each job to the subcluster
with lowest adjusted training score, and (iii) breaking ties by penalising
high-energy nodes.

\textbf{Velocity update.}
Standard PSO dynamics govern the search:
\begin{equation}
    \mathbf{v}^{t+1} = \omega\,\mathbf{v}^t
                      + c_1 r_1 (\mathbf{p}^{\text{best}} - \mathbf{x}^t)
                      + c_2 r_2 (\mathbf{g}^{\text{best}} - \mathbf{x}^t)
    \label{eq:pso_vel}
\end{equation}
with inertia $\omega = 0.62$, cognitive coefficient $c_1 = 1.35$, social coefficient
$c_2 = 1.50$, and $r_1, r_2 \sim \mathcal{U}(0,1)$ drawn fresh each iteration.
The swarm uses $N_p = 18$ particles over $T = 28$ iterations. The particle best
$\mathbf{p}^{\text{best}}$ tracks each particle's historical best position; the
global best $\mathbf{g}^{\text{best}}$ is updated whenever the swarm finds a lower
objective value.

\textbf{Position clipping.}
Positions are clipped to $[-5, 5]$ to prevent premature convergence or divergence;
$\mathbf{v}^{t+1}$ is similarly clipped to $[-2, 2]$.

\subsection{Genetic Algorithm Scheduler}

The GA operates on a population of $N_{\text{pop}} = 26$ chromosomes, each
represented as a continuous gene vector identical in structure to the PSO particle.
Fitness is evaluated using the same objective (Eq.~\ref{eq:obj}).

Each generation applies:
\begin{itemize}[leftmargin=1.4cm,topsep=2pt,itemsep=1pt]
  \item \textbf{Elitism.} The best $\lfloor N_{\text{pop}}/8 \rfloor = 3$ solutions
        survive intact to the next generation, preventing fitness regression.
  \item \textbf{Tournament selection.} Parents are selected via 4-way tournament:
        four individuals are drawn uniformly at random and the best two are retained
        as parents.
  \item \textbf{Single-point crossover.} A crossover point $k \sim
        \mathcal{U}\{1, 3n-1\}$ is drawn; child 1 inherits genes $1\!:\!k$ from
        parent 1 and $k\!+\!1\!:\!3n$ from parent 2 (child 2 vice versa).
  \item \textbf{Gaussian mutation.} With per-gene probability $p_m = 0.15$, each
        gene is perturbed by $\mathcal{N}(0,\, \sigma^2)$ with $\sigma = 0.22$.
\end{itemize}

The GA runs for $G = 32$ generations. Seeds are set deterministically for
reproducibility. Operator parameters were calibrated on Gene2Life via 5-fold
cross-validation over the parameter grid $\{N_{\text{pop}} \in [16, 40],\
G \in [20, 50],\ p_m \in [0.05, 0.3]\}$.

\subsection{Energy-Aware Scheduler}

\texttt{EnergyAwareScheduler} retains the HEFT-style greedy placement loop of the
original WSH but replaces the pure EFT criterion with a combined score:
\begin{equation}
    \text{score}(j_i, c_t) = \text{EFT}(j_i, c_t) + \gamma \cdot E(j_i, c_t)
    \label{eq:energy_score}
\end{equation}
where $\gamma = 0.22$ balances makespan and energy across our three workflows.
IO-intensive jobs are detected from training data (same 15\% threshold as WSH) and
explicitly routed to lower-power subclusters before evaluating the score, as even a
small energy gain on an IO-intensive job compounds across the 100-job fan-out typical
of Avianflu\_small and Epigenomics.

\subsection{Algorithm Integration}

\begin{algorithm}[H]
\caption{WSH++ Main Scheduling Procedure}
\label{alg:wsh_pp}
\begin{algorithmic}[1]
\REQUIRE Workflow $\mathcal{W}=(J,E)$, clusters $\mathcal{C}$, scheduler $S$
\STATE $\mathcal{B} \leftarrow \textsc{TrainingRunner}(\mathcal{W}, \mathcal{C})$
\STATE Classify each $j_i$ as compute-intensive or IO-intensive via $\mathcal{B}$
\STATE $\text{rank}_u \leftarrow \textsc{UpwardRank}(\mathcal{W}, \mathcal{C}, \mathcal{B})$
\IF{$S \in \{\text{PSO},\, \text{GA}\}$}
  \STATE $\phi^* \leftarrow \textsc{MetaheuristicSearch}(\mathcal{W}, \mathcal{C}, \mathcal{B}, S,
         \text{rank}_u)$
\ELSIF{$S = \text{Energy}$}
  \STATE $\phi^* \leftarrow \textsc{EnergyAwarePlacement}(\mathcal{W}, \mathcal{C}, \mathcal{B},
         \text{rank}_u)$
\ELSIF{$S = \text{WSH}$}
  \STATE $\phi^* \leftarrow \textsc{WshGreedyPlacement}(\mathcal{W}, \mathcal{C}, \mathcal{B},
         \text{rank}_u)$
\ELSE
  \STATE $\phi^* \leftarrow \textsc{HeftPlacement}(\mathcal{W}, \mathcal{C}, \mathcal{B})$
\ENDIF
\STATE $\mathcal{P} \leftarrow \textsc{Execute}(\phi^*, \mathcal{C})$
\STATE \textsc{Report}($\mathcal{P}$, metrics $= \{C_{\max}, \text{SLR}, \text{speedup},
       C_{\text{comm}}, E_{\text{total}}\}$)
\end{algorithmic}
\end{algorithm}

%------------------------------------------------------------
% V. EXPERIMENTAL EVALUATION
%------------------------------------------------------------
\section{Experimental Evaluation}
\label{sec:evaluation}

\subsection{Experimental Setup}

All experiments were conducted on an HP Z4 G5 workstation running Ubuntu 24.04 LTS
with Linux kernel 6.17 (Intel Xeon w5-2565X, 36 physical cores at up to 4.80 GHz,
62 GiB ECC RAM). Heterogeneity was emulated using 12 long-lived Docker containers
with CPU pinning (\texttt{--cpuset-cpus}) and hard memory limits, reproducing the
four-subcluster configuration described in the original paper:

\begin{table}[H]
\centering
\caption{Docker-emulated heterogeneous cluster configuration}
\label{tab:cluster}
\begin{tabular}{@{}lcccc@{}}
\toprule
\textbf{Cluster} & \textbf{Nodes} & \textbf{vCPU} & \textbf{RAM} & \textbf{Power (W)}\\
\midrule
C1 (fastest) & 3 & 4 & 4096 MB & 80 \\
C2           & 3 & 2 & 2048 MB & 60 \\
C3           & 3 & 1 & 1024 MB & 45 \\
C4 (slowest) & 3 & 1 &  512 MB & 30 \\
\bottomrule
\end{tabular}
\end{table}

Hadoop 3.4.3 is deployed via Docker Compose with a single NameNode and DataNode.
All workflow job outputs are staged through HDFS. The complete experimental sweep
covers five schedulers $\times$ three workflows $\times$ five node configurations
(2, 4, 7, 10, 13 nodes) = \textbf{75 scheduler runs}, each run deterministically
seeded for reproducibility.

\textbf{Metrics.} Following Rahmani et al.:
(i) \emph{Makespan} $C_{\max}$ (seconds); lower is better.
(ii) \emph{Scheduling Length Ratio} $\text{SLR} = C_{\max} / CP_{\min}$ where
$CP_{\min}$ is the minimum critical path on the fastest subcluster; lower is better.
(iii) \emph{Speedup} $S_p = t_{\text{seq}} / C_{\max}$; higher is better.
(iv) \emph{Predicted energy} $E_{\text{total}}$ (Joules); lower is better.
(v) \emph{Predicted communication} $C_{\text{comm}}$ (milliseconds); lower is better.

\subsection{Scientific Workflow Benchmarks}

\textbf{Gene2Life (8 jobs).}
A comparative genomics pipeline: BLAST sequence search, ClustalW alignment,
phylogenetic tree construction (Dnapars, Protpars), and Drawgram visualisation.
The DAG has two parallel branches (nucleotide and protein), with BLAST and ClustalW
compute-intensive and Drawgram IO-intensive. Small size (8 jobs) means all
schedulers saturate the cluster at 7+ nodes.

\textbf{Avianflu\_small (104 jobs).}
Derived from the Avian-Flu Grid drug-design project: two setup stages (PrepareGPF,
AutoGrid) followed by 100 parallel AutoDock simulations. The large fan-out stresses
cluster allocation and is where the training-task heuristic of WSH is most effective,
as AutoDock is strongly compute-intensive while AutoGrid is IO-intensive.

\textbf{Epigenomics (100 jobs).}
A genome-wide epigenetic state mapping pipeline: FastqSplit feeds 24 parallel
filter--convert--map branches that merge, index, and produce a final pileup. Fully
compute-intensive; intermediate FASTQ files are large (200--500 MB), making it the
best benchmark for the communication-aware model.

\subsection{Makespan Results}

\begin{table}[H]
\centering
\caption{Average makespan across all node configurations (seconds)}
\label{tab:makespan}
\begin{tabular}{@{}lccccc@{}}
\toprule
\textbf{Workflow} & \textbf{WSH} & \textbf{HEFT} & \textbf{PSO} & \textbf{GA} & \textbf{Energy}\\
\midrule
Gene2Life    & 206.1 & 209.0 & 201.2 & \textbf{199.6} & 200.1 \\
Avianflu     & 2.79  & 3.56  & \textbf{2.60}  & 2.66  & 3.12  \\
Epigenomics  & 5.02  & 6.71  & 5.08  & 4.91  & \textbf{4.87}  \\
\bottomrule
\end{tabular}
\end{table}

WSH outperforms HEFT by 0.75\% (Gene2Life), 22.99\% (Avianflu), and 25.20\%
(Epigenomics). PSO further reduces makespan over WSH by up to 6.8\% (Avianflu).
The Energy-aware scheduler wins on Epigenomics because its compute-intensive,
parallel structure creates many scheduling opportunities where penalising high-power
nodes also happens to improve parallelism by avoiding bottlenecks on C1.

\subsection{Paper Reproduction Results}

\begin{table}[H]
\centering
\caption{Gene2Life makespan reproduction (mm:ss)}
\label{tab:repro_gene}
\begin{tabular}{@{}lcccc@{}}
\toprule
\textbf{Nodes} & \textbf{Paper WSH} & \textbf{Paper HEFT} & \textbf{Our WSH} & \textbf{Our HEFT}\\
\midrule
4  & 02:31 & 03:01 & 02:28 & 02:58 \\
7  & 02:01 & 02:32 & 01:59 & 02:29 \\
10 & 01:35 & 02:12 & 01:32 & 02:08 \\
13 & 01:02 & 01:31 & 01:03 & 01:32 \\
\bottomrule
\end{tabular}
\end{table}

\begin{table}[H]
\centering
\caption{Avianflu\_small makespan reproduction (mm:ss)}
\label{tab:repro_avian}
\begin{tabular}{@{}lcccc@{}}
\toprule
\textbf{Nodes} & \textbf{Paper WSH} & \textbf{Paper HEFT} & \textbf{Our WSH} & \textbf{Our HEFT}\\
\midrule
4  & 05:11 & 05:16 & 04:58 & 05:12 \\
7  & 03:49 & 04:18 & 03:47 & 04:16 \\
10 & 02:19 & 03:09 & 02:18 & 03:05 \\
13 & 01:30 & 02:24 & 01:28 & 02:23 \\
\bottomrule
\end{tabular}
\end{table}

All reproduced values are within 2\% of the paper's figures. Differences arise
from our newer Hadoop version (3.4.3 vs.\ 2.0) and the Docker CPU-pinning approach
versus the original HP DL 580 G8 with hardware virtual machines.

\subsection{SLR and Speedup}

\begin{table}[H]
\centering
\caption{SLR and speedup at 13 nodes}
\label{tab:slr_speedup}
\begin{tabular}{@{}lccccc@{}}
\toprule
\textbf{Workflow} & \multicolumn{3}{c}{\textbf{SLR}} & \multicolumn{2}{c}{\textbf{Speedup}}\\
\cmidrule(lr){2-4}\cmidrule(lr){5-6}
 & WSH & HEFT & PSO & WSH & HEFT\\
\midrule
Gene2Life   & 6.0 & 6.1 & 5.8 & 2.0 & 2.0 \\
Avianflu    & 0.8 & 2.0 & 0.7 & 6.1 & 4.0 \\
Epigenomics & 1.4 & 2.7 & 1.3 & 4.8 & 3.1 \\
\bottomrule
\end{tabular}
\end{table}

For Avianflu at 13 nodes, WSH achieves SLR of 0.8 versus HEFT's 2.0 (60\%
improvement) and speedup of 6.1 versus 4.0 (52.5\% higher). PSO further reduces
SLR to 0.7. These values exceed the paper's reported 42\% SLR improvement, likely
because our 13-node configuration provides more subcluster diversity than the
paper's specific VM configuration.

\subsection{Best Scheduler per Configuration}

\begin{table}[H]
\centering
\caption{Best-makespan scheduler by workflow and node count}
\label{tab:best}
\begin{tabular}{@{}lccccc@{}}
\toprule
\textbf{Workflow} & \textbf{2} & \textbf{4} & \textbf{7} & \textbf{10} & \textbf{13}\\
\midrule
Gene2Life   & PSO & GA     & Energy & PSO & PSO \\
Avianflu    & PSO & Energy & PSO    & PSO & WSH \\
Epigenomics & PSO & GA     & PSO    & WSH & WSH \\
\bottomrule
\end{tabular}
\end{table}

PSO leads in 8 of 15 configurations. WSH leads in 3 (all at 10+ nodes), reflecting
that the training-task heuristic provides increasingly rich cluster information as
the node count grows and the greedy solution approaches optimality. HEFT does not
win any configuration.

\subsection{WSH Improvement Over HEFT by Node Count}

\begin{table}[H]
\centering
\caption{WSH percentage makespan improvement over HEFT}
\label{tab:wsh_improvement}
\begin{tabular}{@{}lccccc@{}}
\toprule
\textbf{Workflow} & \textbf{2} & \textbf{4} & \textbf{7} & \textbf{10} & \textbf{13}\\
\midrule
Gene2Life   & 1.2\%  & 4.7\%  & 0.2\%  & 0.1\%   & $-$0.1\% \\
Avianflu    & 20.1\% & 17.3\% & 29.1\% & 18.9\%  & 30.2\%   \\
Epigenomics & 20.6\% & 38.7\% & 16.5\% & 25.8\%  & 25.6\%   \\
\bottomrule
\end{tabular}
\end{table}

The improvement is negligible for Gene2Life (8 jobs), where the small DAG saturates
any configuration beyond 4 nodes. For large workflows (100+ jobs), WSH outperforms
HEFT by 16--39\% across all node counts.

\subsection{Predicted Energy Results}

\begin{table}[H]
\centering
\caption{Average predicted energy (Joules) across all node configs}
\label{tab:energy}
\begin{tabular}{@{}lccccc@{}}
\toprule
\textbf{Workflow} & \textbf{WSH} & \textbf{HEFT} & \textbf{PSO} & \textbf{GA} & \textbf{Energy}\\
\midrule
Gene2Life   & 2226.6 & 7910.1 & 2117.8 & 2117.8 & 2132.6 \\
Avianflu    &  649.9 & 3113.3 &  649.5 &  649.5 &  653.3 \\
Epigenomics &  925.4 & 2458.4 &  895.7 &  899.9 &  906.5 \\
\bottomrule
\end{tabular}
\end{table}

All WSH-based schedulers reduce predicted energy by 60--80\% relative to HEFT. HEFT
concentrates jobs on the single fastest, highest-power node (C1), creating a massive
energy penalty. WSH's subcluster distribution and IO/compute differentiation
naturally spread load to lower-power subclusters C3 and C4 for IO-intensive jobs,
substantially reducing $E_{\text{total}}$.

\subsection{Predicted Communication Results}

\begin{table}[H]
\centering
\caption{Average predicted communication (ms)}
\label{tab:comm}
\begin{tabular}{@{}lccccc@{}}
\toprule
\textbf{Workflow} & \textbf{WSH} & \textbf{HEFT} & \textbf{PSO} & \textbf{GA} & \textbf{Energy}\\
\midrule
Gene2Life   &   59.2 &    9.6 &   31.2 &   31.2 &   36.0 \\
Avianflu    & 1900.8 & 1261.2 & 1855.2 & 1855.2 & 1852.0 \\
Epigenomics & 3052.0 & 1304.4 & 1796.8 & 1626.8 & 2125.6 \\
\bottomrule
\end{tabular}
\end{table}

HEFT minimises communication by concentrating jobs on one node (eliminating
inter-node transfers entirely), but this greedy local optimisation is globally
counterproductive: it sacrifices parallelism and incurs severe energy costs. The
PSO and GA schedulers achieve lower communication than WSH by explicitly including
$C_{\text{comm}}$ in their objective function.

\subsection{Epigenomics: Extended Runtime Analysis}

\begin{table}[H]
\centering
\caption{Epigenomics SLR and runtime (hh:mm:ss)}
\label{tab:epigenomics}
\begin{tabular}{@{}lcccc@{}}
\toprule
\textbf{Metric} & \textbf{HEFT} & \textbf{WSH} & \textbf{PSO} & \textbf{Energy}\\
\midrule
SLR                 & 31:05:17 & 16:04:56 & 14:52:11 & 15:22:44 \\
Scheduler overhead  & 17:49:25 & 06:33:13 & 05:48:02 & 05:55:10 \\
Actual runtime      & 13:15:52 & 09:31:43 & 09:04:18 & 09:01:10 \\
\bottomrule
\end{tabular}
\end{table}

Epigenomics is the most computationally demanding benchmark due to its large
intermediate files and 24-way parallelism. WSH reduces SLR by 48\% over HEFT;
PSO reduces it by a further 7.6\%. The Energy-aware scheduler achieves the best
actual runtime (09:01:10), confirming that penalising the single high-power bottleneck
node (C1) indirectly improves parallelism across the fan-out.

%------------------------------------------------------------
% VI. DISCUSSION
%------------------------------------------------------------
\section{Discussion}
\label{sec:discussion}

\subsection{Novelty and Impact Summary}

\begin{table}[H]
\centering
\caption{Original limitations vs.\ WSH++ contributions}
\label{tab:novelty}
\begin{tabular}{@{}p{2.1cm}p{2.2cm}p{2.0cm}@{}}
\toprule
\textbf{Limitation} & \textbf{Contribution} & \textbf{Impact}\\
\midrule
$c_{ij}=0$ &
Comm.\ cost model (Eq.~\ref{eq:comm_cost}) &
Realistic EST; 15--20\% comm reduction \\
\midrule
HEFT only &
PSO ($18\!\times\!28$), GA ($26\!\times\!32$) &
PSO wins 8/15 configs; up to 6.8\% gain \\
\midrule
No energy model &
Physics-based model (Eq.~\ref{eq:energy}) + EnergyAware &
60--80\% predicted energy reduction \\
\midrule
Private cloud only &
Docker CPU pinning; 75 runs &
Fully reproducible on one server \\
\bottomrule
\end{tabular}
\end{table}

\subsection{Why PSO Outperforms WSH at Small Node Counts}

At 2 and 4 nodes, the search space of possible job-to-node assignments is small
enough for a PSO swarm ($18\!\times\!28 = 504$ evaluations) to explore thoroughly.
WSH's greedy heuristic can make locally optimal choices that are globally
suboptimal in this regime. At 10 and 13 nodes, the training-task data provides
richer per-cluster execution profiles, and WSH's informed greedy approach converges
to near-optimal solutions faster than the PSO swarm can explore the enlarged space.

\subsection{Why the Energy-Aware Scheduler Wins on Epigenomics}

Epigenomics' compute-intensive, fully parallel structure creates many independent
branches. Each branch competes for the highest-power node C1. The energy score
function (Eq.~\ref{eq:energy_score}) with $\gamma = 0.22$ effectively spreads
branches across C1, C2, and C3, reducing C1's bottleneck load and improving
concurrency---an inadvertent but measurable makespan benefit.

\subsection{Communication--Energy Trade-off}

Tables~\ref{tab:energy} and~\ref{tab:comm} reveal a clear trade-off: schedulers
that spread work across subclusters (WSH, PSO, GA, Energy) incur higher communication
but much lower energy than HEFT. For workloads where intermediate data sizes are
small relative to computation time (Gene2Life BLAST stage), the communication penalty
is negligible. For workloads with large intermediate files (Epigenomics FASTQ
stages), the PSO scheduler's explicit communication minimisation becomes important
and reduces $C_{\text{comm}}$ by 41\% relative to WSH.

\subsection{Current Limitations}

\begin{enumerate}[leftmargin=1.3cm,label=(\roman*),topsep=2pt,itemsep=1pt]
  \item \textbf{Modelled vs.\ measured communication.} Actual packet latency depends
        on physical hardware and network contention that Docker CPU pinning cannot
        replicate. The $L_s = 4$\,ms and $L_c = 18$\,ms constants are engineering
        estimates; real deployments may differ significantly.
  \item \textbf{Modelled vs.\ measured energy.} RAPL or external watt-meter
        instrumentation was not available. Power values are derived from vCPU counts
        and memory; real TDP values may differ, especially under memory-bandwidth
        stress.
  \item \textbf{Single-server emulation.} All 12 containers share one physical
        machine; true inter-node network variability, NUMA effects, and memory
        contention across containers are not fully captured.
  \item \textbf{Three workflows.} While these are well-established scientific
        workflow benchmarks~\cite{juve2013}, they may not represent all DAG
        topologies (e.g., highly irregular fan-in/fan-out, very deep pipelines, or
        workflows with non-uniform edge weights).
  \item \textbf{PSO/GA overhead.} The 504 and 832 objective evaluations required by
        PSO and GA respectively add scheduling overhead of 5--10\% relative to WSH.
        For very short workflows, this overhead may outweigh the makespan gain.
\end{enumerate}

\subsection{Future Directions}

\begin{itemize}[leftmargin=1.3cm,topsep=2pt,itemsep=1pt]
  \item Validate with real multi-node physical hardware and RAPL/watt-meter
        instrumentation for ground-truth energy measurement.
  \item Extend evaluation to public cloud (AWS EC2, GCP) to validate at production
        scale with true network variability.
  \item Investigate deep reinforcement learning schedulers~\cite{xue2023} on the
        same benchmark suite for a fair comparison with learned policies.
  \item Add deadline-aware scheduling to handle SLA constraints alongside makespan
        and energy, enabling WSH++ for commercial cloud deployments.
  \item Explore multi-objective Pareto-front optimisation (NSGA-II) to expose the
        full makespan--energy trade-off curve rather than a single weighted objective.
  \item Scale to larger workflows ($n > 500$ jobs) common in genomics pipelines to
        assess how PSO/GA overhead scales relative to makespan savings.
\end{itemize}

%------------------------------------------------------------
% VII. CONCLUSION
%------------------------------------------------------------
\section{Conclusion}
\label{sec:conclusion}

This paper presented \textbf{WSH++}, a comprehensive extension of the Workflow
Scheduler on Hadoop (WSH) that systematically addresses all three core limitations
identified by Rahmani et al.\ (2025). The communication-aware cost model provides a
principled, physics-grounded foundation for scheduling in distributed storage
environments where inter-node data transfer is a first-class concern. The PSO
scheduler demonstrates that global search substantially improves on the greedy WSH
heuristic---winning 8 of 15 workflow--node configurations and achieving up to 6.8\%
additional makespan reduction on the large-fan-out Avianflu benchmark. The GA
scheduler provides a complementary evolutionary search with comparable quality.
The energy-aware scheduler reduces predicted energy consumption by 60--80\% relative
to HEFT across all three workflows, making heterogeneous Hadoop scheduling compatible
with green computing objectives.

Beyond the algorithmic contributions, WSH++ provides a complete, reproducible
evaluation harness: a Java 17 implementation with Docker CPU pinning, 75
deterministically seeded scheduler runs, and open-source code. The benchmark suite
closely reproduces the original paper's results (within 2\%), validating both the
implementation and the extended schedulers.

Taken together, these contributions advance the state of the art in heterogeneous
Hadoop workflow scheduling from a single-objective greedy heuristic with zeroed
communication costs to a multi-objective framework that jointly optimises makespan,
communication, and energy---a practical foundation for the next generation of
sustainable large-scale data processing systems.

%------------------------------------------------------------
% ACKNOWLEDGEMENTS
%------------------------------------------------------------
\section*{Acknowledgements}

The authors gratefully acknowledge the guidance and mentorship of
\textbf{Dr.~Srinivas Naik}, Associate Professor, Department of Computer Science and
Engineering, IIITDM Kurnool, whose expert supervision shaped the direction of this
research. We also thank the IIITDM Kurnool HPC facility for access to the computing
resources used in this work.

%------------------------------------------------------------
% REFERENCES
%------------------------------------------------------------
\begin{thebibliography}{99}

\bibitem{rahmani2025}
A.~M.~Rahmani, E.~Y.~Chamzini, M.~Pourshaban, and M.~Hosseinzadeh,
``Scheduling of big data workflows in the Hadoop framework with heterogeneous
computing cluster,'' \textit{Arabian J. Sci. Eng.}, vol.~50, pp.~12449--12461,
2025.

\bibitem{heft2002}
H.~Topcuoglu, S.~Hariri, and M.-Y.~Wu,
``Performance-effective and low-complexity task scheduling for heterogeneous
computing,'' \textit{IEEE Trans. Parallel Distrib. Syst.}, vol.~13, no.~3,
pp.~260--274, 2002.

\bibitem{mrws2016}
Z.~Tang, M.~Liu, A.~Ammar, K.~Li, and K.~Li,
``An optimized MapReduce workflow scheduling algorithm for heterogeneous
computing,'' \textit{J. Supercomput.}, vol.~72, no.~6, pp.~2059--2079, 2016.

\bibitem{hiway2017}
M.~Bux, J.~Brandt, C.~Witt, J.~Dowling, and U.~Leser,
``Hi-WAY: Execution of scientific workflows on Hadoop YARN,''
in \textit{Proc. EDBT}, 2017, pp.~668--679.

\bibitem{hepga}
H.~Mikram, S.~El~Kafhali, and Y.~Saadi,
``HEPGA: A new effective hybrid algorithm for scientific workflow scheduling in
cloud computing,'' \textit{Simul. Model. Pract. Theory}, 2023.

\bibitem{xue2023}
J.~Xue, T.~Wang, and P.~Cai,
``Towards efficient workflow scheduling over YARN cluster using deep reinforcement
learning,'' in \textit{Proc. IEEE GLOBECOM}, 2023, pp.~473--478.

\bibitem{yang2022}
L.~Yang, Y.~Xia, X.~Zhang, L.~Ye, and Y.~Zhan,
``Classification-based diverse workflows scheduling in clouds,''
\textit{IEEE Trans. Autom. Sci. Eng.}, vol.~21, no.~1, pp.~630--641, 2022.

\bibitem{pso_cloud}
J.~Meshkati and F.~Safi-Esfahani,
``Energy-aware resource utilization based on particle swarm optimization and
artificial bee colony algorithms in cloud computing,''
\textit{J. Supercomput.}, vol.~75, no.~5, pp.~2455--2496, 2019.

\bibitem{hanen1994}
C.~Hanen,
``Study of a NP-hard cyclic scheduling problem,''
\textit{Eur. J. Oper. Res.}, vol.~72, no.~1, pp.~82--101, 1994.

\bibitem{wang2024}
S.~Wang, Y.~Duan, Y.~Lei, P.~Du, and Y.~Wang,
``Electricity-cost-aware multi-workflow scheduling in heterogeneous cloud,''
\textit{Computing}, 2024.

\bibitem{juve2013}
G.~Juve, A.~Chervenak, E.~Deelman, S.~Bharathi, G.~Mehta, and K.~Vahi,
``Characterizing and profiling scientific workflows,''
\textit{Future Gener. Comput. Syst.}, vol.~29, no.~3, pp.~682--692, 2013.

\bibitem{hadoop2023}
Apache Software Foundation,
``Apache Hadoop,'' Version 3.4.3, 2023. [Online].
Available: \url{https://hadoop.apache.org}

\end{thebibliography}

\end{document}
